\chapter{Contexte et but du projet}

% ----------------------------------------------------------------
% Introduction à la section
% ----------------------------------------------------------------
% Cette section vise à présenter de manière complète le cadre dans lequel s’inscrit votre projet.
% Elle permet de bien comprendre pourquoi le projet a été proposé, ce qu’il vise à accomplir,
% et dans quel contexte il a émergé (technique, académique, ou professionnel).
%
% Vous pouvez ici résumer en quelques lignes ce qui sera détaillé dans les sous-sections.

% Exemple d'introduction :
% Le projet présenté dans ce document a émergé d’un besoin identifié au sein du campus de Lyon :
% disposer d’un outil pédagogique permettant aux étudiants d’expérimenter concrètement les principes de la cybersécurité
% via l’analyse d’une application existante. Cette section vise à poser les bases de ce projet,
% à travers la description du contexte, de ses objectifs, et des ressources disponibles.

% ----------------------------------------------------------------
% Import des sous-sections
% ----------------------------------------------------------------
\section{Objectif du projet}

% ----------------------------------------------------------------
% 📌 À quoi sert cette section ?
% ----------------------------------------------------------------
% Cette section doit répondre clairement à la question :
% « Que souhaite-t-on accomplir avec ce projet ? »
% 
% Il ne s’agit pas encore de détailler les fonctionnalités ni les solutions techniques, 
% mais de donner une vision synthétique et globale des finalités du projet.
%
% L’objectif doit être SMART (Spécifique, Mesurable, Atteignable, Réaliste, Temporellement défini),
% ou du moins respecter ces grandes lignes pour aider le lecteur à comprendre l’utilité du projet.

% ----------------------------------------------------------------
% 🧭 Comment la remplir ?
% ----------------------------------------------------------------
% ➤ Reformulez l’objectif du projet en une ou deux phrases d’introduction.
% ➤ Ajoutez ensuite un développement précisant :
%     - Le public cible
%     - Le contexte d’utilisation
%     - Le besoin ou la motivation ayant mené à ce projet
%     - Les grandes lignes du résultat attendu
%
% ➤ Vous pouvez, si cela aide à clarifier, ajouter un schéma illustrant le positionnement
%   ou la finalité du projet dans son environnement.

% ----------------------------------------------------------------
% 🧾 Exemple rédigé :
% ----------------------------------------------------------------
% L’objectif du projet est de concevoir un outil interactif permettant aux étudiants de s’initier
% au langage SVG à travers une interface ludique de type "bac à sable".
% Le but est de rendre l’apprentissage de la syntaxe SVG plus intuitif, tout en conservant
% une liberté de création.
%
% Ce projet s’adresse aux étudiants de première année en informatique, dans un contexte éducatif.
% Le livrable visé est une plateforme accessible en ligne, ne nécessitant aucune installation locale,
% et intégrant des défis progressifs ainsi qu’une prévisualisation en temps réel.

% ----------------------------------------------------------------
% 📊 Exemple de schéma à insérer (schéma positionnement)
% ----------------------------------------------------------------

% \begin{figure}[H]
%     \centering
%     \begin{tikzpicture}[node distance=1.5cm]
%         \node (start) [draw, rectangle] {Problème identifié : apprentissage SVG difficile};
%         \node (middle) [below of=start, draw, rectangle] {Solution proposée : plateforme interactive};
%         \node (end) [below of=middle, draw, rectangle] {Objectif final : apprentissage ludique et progressif};
%         \draw[->] (start) -- (middle);
%         \draw[->] (middle) -- (end);
%     \end{tikzpicture}
%     \caption{Schéma simplifié de l’objectif du projet}
%     \label{fig:objectif}
% \end{figure}

% ----------------------------------------------------------------
% ✍️ À vous de jouer : complétez ci-dessous
% ----------------------------------------------------------------

% (Exemple de rédaction)
L’objectif du projet est de...


\section{À quel problème répond votre projet ?}

% ----------------------------------------------------------------
% 📌 À quoi sert cette section ?
% ----------------------------------------------------------------
% Cette partie permet de contextualiser l'objectif du projet en expliquant :
% - Le problème initial qui a motivé sa création
% - Les limites ou lacunes des solutions existantes
% - Les enjeux techniques, humains ou pédagogiques concernés
%
% C’est une section essentielle qui justifie l’intérêt du projet.

% ----------------------------------------------------------------
% 🧭 Comment la remplir ?
% ----------------------------------------------------------------
% ➤ Identifiez clairement un problème ou une frustration vécue par les utilisateurs ou les acteurs concernés.
% ➤ Expliquez en quoi ce problème est réel, fréquent ou bloquant.
% ➤ Si possible, illustrez avec des cas d’usage, des témoignages, ou des situations observées.
% ➤ Terminez par une phrase de transition vers la solution proposée dans le projet.

% ----------------------------------------------------------------
% 🧾 Exemple rédigé :
% ----------------------------------------------------------------
% Aujourd’hui, l’apprentissage du format SVG reste très théorique dans les cursus de développement front-end.
% Les ressources disponibles sont nombreuses mais peu interactives et rarement adaptées aux débutants.
% Il n’existe pas de plateforme ludique francophone permettant d’expérimenter directement avec le code SVG
% sans configuration complexe.
%
% Ce manque de support pratique freine l’autonomisation des étudiants et limite leur motivation.
% Le projet vise à combler cette lacune avec un outil visuel accessible en ligne.

% ----------------------------------------------------------------
% 📊 Exemple de schéma à insérer (problème & causes)
% ----------------------------------------------------------------

% \begin{figure}[H]
%     \centering
%     \begin{tikzpicture}[node distance=1.5cm]
%         \node (pb) [draw, rectangle, fill=red!10] {Problème : apprentissage SVG non pratique};
%         \node (c1) [below left=of pb, draw, rectangle] {Absence d'outils interactifs};
%         \node (c2) [below right=of pb, draw, rectangle] {Ressources trop complexes};
%         \draw[->] (c1) -- (pb);
%         \draw[->] (c2) -- (pb);
%     \end{tikzpicture}
%     \caption{Causes du problème identifié}
%     \label{fig:problematique}
% \end{figure}

% ----------------------------------------------------------------
% ✍️ À vous de jouer : complétez ci-dessous
% ----------------------------------------------------------------

% (Exemple de départ)
Le problème identifié dans le cadre de ce projet est que...


\section{Documentation de l'existant}

% ----------------------------------------------------------------
% 📌 À quoi sert cette section ?
% ----------------------------------------------------------------
% Cette partie sert à prouver que vous avez exploré les solutions déjà disponibles
% (outils, logiciels, méthodes, articles) en lien avec le problème que vous abordez.
% Elle permet de montrer que votre projet est justifié par un manque, une faiblesse
% ou une opportunité d'amélioration dans ce qui existe déjà.

% ----------------------------------------------------------------
% 🧭 Comment la remplir ?
% ----------------------------------------------------------------
% ➤ Listez les outils ou approches existants en lien avec votre projet.
% ➤ Décrivez leurs fonctionnalités principales.
% ➤ Soulignez leurs avantages ET leurs limites dans le cadre de votre problématique.
% ➤ Comparez-les entre eux si pertinent.
% ➤ Concluez par l’intérêt de votre projet comme complément, alternative ou amélioration.

% Vous pouvez aussi ajouter un tableau comparatif ou un schéma d’analyse.

% ----------------------------------------------------------------
% 🧾 Exemple rédigé :
% ----------------------------------------------------------------
% Plusieurs plateformes de test en ligne existent pour manipuler du code SVG :
% - \textbf{CodePen} : puissant et flexible, mais non spécialisé SVG
% - \textbf{SVGPlayground} : permet de tester du SVG, mais sans objectif éducatif
% - \textbf{JSFiddle} : généraliste, sans guidage pédagogique
%
% Ces outils ont en commun leur accessibilité, mais manquent d’accompagnement pour des profils débutants.
% Aucun ne propose une approche gamifiée ou structurée autour de défis pédagogiques.

% ----------------------------------------------------------------
% 📊 Exemple de tableau comparatif
% ----------------------------------------------------------------
% \begin{table}[H]
% \centering
% \begin{tabular}{|l|c|c|c|}
% \hline
% Outil & Ciblé débutant & Pédagogique & Temps réel \\
% \hline
% CodePen & Non & Non & Oui \\
% SVGPlayground & Partiellement & Non & Oui \\
% Projet proposé & Oui & Oui & Oui \\
% \hline
% \end{tabular}
% \caption{Comparaison des solutions existantes}
% \end{table}

% ----------------------------------------------------------------
% ✍️ À vous de jouer : complétez ci-dessous
% ----------------------------------------------------------------

% (Début de réponse)
À ce jour, plusieurs outils existent dans le domaine, parmi lesquels...


\section{Quels sont les apports de ce projet ?}

% ----------------------------------------------------------------
% 📌 À quoi sert cette section ?
% ----------------------------------------------------------------
% Cette section met en avant la valeur ajoutée du projet :
% - En quoi ce projet est-il bénéfique pour les étudiants porteurs ?
% - Quelles compétences ou connaissances seront développées ?
% - Quel est l’apport pour l’école, pour une entreprise, ou pour un utilisateur final ?
%
% Elle peut aborder des aspects techniques, pédagogiques, humains ou organisationnels.

% ----------------------------------------------------------------
% 🧭 Comment la remplir ?
% ----------------------------------------------------------------
% ➤ Listez les compétences que vous allez mobiliser ou acquérir (techniques et non techniques).
% ➤ Mentionnez les apports pour votre profil (développement, gestion, UX, communication...).
% ➤ Évoquez la valeur ajoutée pour SUPINFO, un partenaire, ou une communauté cible.
% ➤ Vous pouvez illustrer l’évolution ou l’enrichissement du profil de l’étudiant avec un tableau simple ou un schéma.

% ----------------------------------------------------------------
% 🧾 Exemple rédigé :
% ----------------------------------------------------------------
% Ce projet permettra aux étudiants de développer plusieurs compétences techniques, telles que :
% - La maîtrise de React pour la construction d’interfaces web interactives
% - L’utilisation de bibliothèques telles que DOMPurify pour sécuriser du contenu dynamique
% - La capacité à concevoir une architecture pédagogique structurée
%
% En parallèle, des compétences transversales seront renforcées :
% - Communication technique via la rédaction de documentation
% - Gestion de projet en mode agile (backlog, user stories)
% - Collaboration avec des enseignants ou partenaires externes
%
% Ce projet peut également devenir un outil pédagogique réutilisable à SUPINFO pour les promotions futures.

% ----------------------------------------------------------------
% 📊 Exemple de tableau d'apports (techniques / transversaux)
% ----------------------------------------------------------------

% \begin{table}[H]
% \centering
% \begin{tabular}{|l|p{8cm}|}
% \hline
% Domaine & Compétences développées \\
% \hline
% Technique & Programmation React, gestion des dépendances JS, sandbox sécurisé \\
% Pédagogique & Structuration de parcours d’apprentissage, gamification \\
% Transversal & Rédaction technique, gestion du temps, collaboration \\
% \hline
% \end{tabular}
% \caption{Synthèse des apports du projet}
% \end{table}

% ----------------------------------------------------------------
% ✍️ À vous de jouer : complétez ci-dessous
% ----------------------------------------------------------------

% (Début de réponse)
Ce projet permet aux porteurs de développer...


\section{Porteurs du projet}

% ----------------------------------------------------------------
% 📌 À quoi sert cette section ?
% ----------------------------------------------------------------
% Cette partie sert à identifier formellement les étudiants responsables du projet.
% Elle permet de les associer clairement à la production, pour le suivi et la validation.

% ----------------------------------------------------------------
% 🧭 Comment la remplir ?
% ----------------------------------------------------------------
% ➤ Indiquez les informations suivantes pour chaque porteur :
%     - Prénom NOM
%     - Promotion (ex: M1, M2)
%     - Campus SUPINFO
%     - Adresse e-mail institutionnelle (prenom.nom@supinfo.com)

% ➤ Si vous êtes plusieurs, vous pouvez organiser l’information sous forme de liste ou de tableau.

% ----------------------------------------------------------------
% 🧾 Exemple sous forme de liste :
% ----------------------------------------------------------------
% \begin{itemize}
%   \item Rémi COVIZZI – M2 – Lyon – remi.covizzi@supinfo.com
%   \item Sophie MARTIN – M1 – Paris – sophie.martin@supinfo.com
% \end{itemize}

% ----------------------------------------------------------------
% 📊 Exemple sous forme de tableau :
% ----------------------------------------------------------------
% \begin{table}[H]
% \centering
% \begin{tabular}{|l|c|c|l|}
% \hline
% Nom & Promotion & Campus & Email \\
% \hline
% Rémi COVIZZI & M2 & Lyon & remi.covizzi@supinfo.com \\
% Sophie MARTIN & M1 & Paris & sophie.martin@supinfo.com \\
% \hline
% \end{tabular}
% \caption{Étudiants porteurs du projet}
% \end{table}

% ----------------------------------------------------------------
% ✍️ À vous de jouer : complétez ci-dessous
% ----------------------------------------------------------------

% (Modèle recommandé)
\begin{itemize}
  \item [À compléter]
\end{itemize}

\section{Objectifs fonctionnels}

% ----------------------------------------------------------------
% 📌 À quoi sert cette section ?
% ----------------------------------------------------------------
% Cette partie liste de manière claire et synthétique les principales fonctionnalités attendues
% ou actions que le système/projet doit permettre à l’utilisateur d’accomplir.
% Il s’agit d’une vue fonctionnelle globale (le "quoi"), sans entrer dans les détails techniques (le "comment").

% ----------------------------------------------------------------
% 🧭 Comment la remplir ?
% ----------------------------------------------------------------
% ➤ Rédigez les objectifs fonctionnels sous forme de liste ou d’User Stories.
% ➤ Les User Stories suivent souvent la forme : 
%     "En tant que [rôle], je veux [fonctionnalité] afin de [bénéfice]."
% ➤ Classez-les par ordre logique ou de priorité si possible.
% ➤ Soyez clairs, concis, et évitez les termes trop techniques.

% ----------------------------------------------------------------
% 🧾 Exemple sous forme d’User Stories :
% ----------------------------------------------------------------
% \begin{itemize}
%   \item En tant qu’étudiant, je veux visualiser le résultat de mon code SVG en temps réel pour mieux comprendre mes erreurs.
%   \item En tant qu’utilisateur débutant, je veux accéder à une documentation simplifiée et interactive.
%   \item En tant que formateur, je veux pouvoir créer des défis personnalisés pour mes élèves.
% \end{itemize}

% ----------------------------------------------------------------
% 📊 Exemple avec tableau de priorités
% ----------------------------------------------------------------
% \begin{table}[H]
% \centering
% \begin{tabular}{|p{9cm}|c|}
% \hline
% Fonctionnalité (User Story) & Priorité \\
% \hline
% Visualiser du code SVG en temps réel & Haute \\
% Intégration de défis gamifiés & Moyenne \\
% Export des résultats & Basse \\
% \hline
% \end{tabular}
% \caption{Objectifs fonctionnels du projet}
% \end{table}

% ----------------------------------------------------------------
% ✍️ À vous de jouer : complétez ci-dessous
% ----------------------------------------------------------------

% (Départ suggéré)
Les objectifs fonctionnels du projet sont les suivants :

\begin{itemize}
  \item [À compléter]
\end{itemize}

\section{Environnement technique}

% ----------------------------------------------------------------
% 📌 À quoi sert cette section ?
% ----------------------------------------------------------------
% Cette partie décrit l’ensemble des outils, technologies, langages et environnements
% utilisés pour mener à bien le projet.
% Elle permet de cadrer le développement et d’assurer la reproductibilité.

% ----------------------------------------------------------------
% 🧭 Comment la remplir ?
% ----------------------------------------------------------------
% ➤ Listez les outils logiciels (IDE, bibliothèques, frameworks, etc.)
% ➤ Indiquez les langages de programmation, OS, navigateurs ou environnements serveur ciblés.
% ➤ Mentionnez les éventuelles contraintes matérielles ou logicielles.
% ➤ Vous pouvez structurer les infos par catégories (langages, frameworks, outils, etc.)

% ----------------------------------------------------------------
% 🧾 Exemple rédigé :
% ----------------------------------------------------------------
% Le projet sera développé à l’aide des technologies suivantes :
%
% \textbf{Langages et Frameworks} :
% \begin{itemize}
%   \item JavaScript / React pour l’interface web
%   \item HTML / CSS pour la structuration visuelle
%   \item SVG natif manipulé en DOM
% \end{itemize}
%
% \textbf{Outils de développement} :
% \begin{itemize}
%   \item VSCode (éditeur de code)
%   \item GitHub pour le versionning
%   \item React CodeMirror pour l’édition sécurisée du SVG
% \end{itemize}
%
% \textbf{Sécurité / Sanitation} :
% \begin{itemize}
%   \item DOMPurify pour filtrer le code SVG injecté
% \end{itemize}
%
% \textbf{Déploiement} :
% \begin{itemize}
%   \item Hébergement via Vercel
% \end{itemize}

% ----------------------------------------------------------------
% 📊 Exemple de structure par tableau
% ----------------------------------------------------------------
% \begin{table}[H]
% \centering
% \begin{tabular}{|l|l|}
% \hline
% Élément & Détail \\
% \hline
% Langage principal & JavaScript \\
% Framework front-end & React \\
% Librairie de code sécurisé & DOMPurify \\
% Environnement de dev & VSCode + Git \\
% Plateforme de déploiement & Vercel \\
% \hline
% \end{tabular}
% \caption{Résumé de l’environnement technique}
% \end{table}

% ----------------------------------------------------------------
% ✍️ À vous de jouer : complétez ci-dessous
% ----------------------------------------------------------------

% (Suggestion de début)
L’environnement technique du projet repose sur les éléments suivants :

\begin{itemize}
  \item [À compléter...]
\end{itemize}

